\documentclass[10pt,a4paper]{amsart}
\usepackage{amsmath,amssymb,amsthm,graphicx}
\usepackage{verbatim}
\newenvironment{reference}{\comment}{\endcomment}
%\newenvironment{reference}{}{}
\newenvironment{slogan}{\comment}{\endcomment}
\newenvironment{history}{\comment}{\endcomment}

% For commutative diagrams we use Xy-pic
\usepackage[all]{xy}

% We use 2cell for 2-commutative diagrams.
\xyoption{2cell}
% Theorem environments.
%
\theoremstyle{plain}
\newtheorem{theorem}[subsection]{Theorem}
\newtheorem{proposition}[subsection]{Proposition}
\newtheorem{lemma}[subsection]{Lemma}

\theoremstyle{definition}
\newtheorem{definition}[subsection]{Definition}
\newtheorem{example}[subsection]{Example}
\newtheorem{exercise}[subsection]{Exercise}
\newtheorem{situation}[subsection]{Situation}

\theoremstyle{remark}
\newtheorem{remark}[subsection]{Remark}
\newtheorem{remarks}[subsection]{Remarks}

\numberwithin{equation}{subsection}

% Macros
%
\def\lim{\mathop{\mathrm{lim}}\nolimits}
\def\colim{\mathop{\mathrm{colim}}\nolimits}
\def\Spec{\mathop{\mathrm{Spec}}}
\def\Hom{\mathop{\mathrm{Hom}}\nolimits}
\def\Ext{\mathop{\mathrm{Ext}}\nolimits}
\def\SheafHom{\mathop{\mathcal{H}\!\mathit{om}}\nolimits}
\def\SheafExt{\mathop{\mathcal{E}\!\mathit{xt}}\nolimits}
\def\Sch{\mathit{Sch}}
\def\Mor{\mathop{\mathrm{Mor}}\nolimits}
\def\Ob{\mathop{\mathrm{Ob}}\nolimits}
\def\Sh{\mathop{\mathit{Sh}}\nolimits}
\def\NL{\mathop{N\!L}\nolimits}
\def\CH{\mathop{\mathrm{CH}}\nolimits}
\def\proetale{{pro\text{-}\acute{e}tale}}
\def\etale{{\acute{e}tale}}
\def\QCoh{\mathit{QCoh}}
\def\Ker{\mathop{\mathrm{Ker}}}
\def\Im{\mathop{\mathrm{Im}}}
\def\Coker{\mathop{\mathrm{Coker}}}
\def\Coim{\mathop{\mathrm{Coim}}}

% Boxtimes
%
\DeclareMathSymbol{\boxtimes}{\mathbin}{AMSa}{"02}

%
% Macros for moduli stacks/spaces
%
\def\QCohstack{\mathcal{QC}\!\mathit{oh}}
\def\Cohstack{\mathcal{C}\!\mathit{oh}}
\def\Spacesstack{\mathcal{S}\!\mathit{paces}}
\def\Quotfunctor{\mathrm{Quot}}
\def\Hilbfunctor{\mathrm{Hilb}}
\def\Curvesstack{\mathcal{C}\!\mathit{urves}}
\def\Polarizedstack{\mathcal{P}\!\mathit{olarized}}
\def\Complexesstack{\mathcal{C}\!\mathit{omplexes}}
% \Pic is the operator that assigns to X its picard group, usage \Pic(X)
% \Picardstack_{X/B} denotes the Picard stack of X over B
% \Picardfunctor_{X/B} denotes the Picard functor of X over B
\def\Pic{\mathop{\mathrm{Pic}}\nolimits}
\def\Picardstack{\mathcal{P}\!\mathit{ic}}
\def\Picardfunctor{\mathrm{Pic}}
\def\Deformationcategory{\mathcal{D}\!\mathit{ef}}
\usepackage{cleveref}
\crefname{originallemma}{original lemma}{original lemmas}
\Crefname{originallemma}{Original lemma}{Original lemmas}
\makeatletter
\newlabel{lemma-Zariski-topology}{{00E0}{source}{}{sourcepointer-1}{}}
\newlabel{lemma-Zariski-topology@cref}{{[originallemma][0][]00E0}{[1][0][]source}{}{}{}}
\makeatother
\crefname{originallemma}{original lemma}{original lemmas}
\Crefname{originallemma}{Original lemma}{Original lemmas}
\makeatletter
\newlabel{lemma-dominate}{{00IA}{source}{}{sourcepointer-2}{}}
\newlabel{lemma-dominate@cref}{{[originallemma][0][]00IA}{[1][0][]source}{}{}{}}
\makeatother
\title{Existence of valuation rings between a local subring and its fraction field}
\author{The Stacks Project Authors}
\date{stacks-project:algebra.tex:local-snapshot-20261012; GNU FDL 1.2 or later}
\begin{document}
\maketitle
\noindent\textit{Editorial note.} All delivered mathematical text is the author's own, copied from the pinned Stacks Project chapter file algebra.tex; no statement or proof step was rewritten, reordered or removed.\par\medskip
\noindent\textit{Editorial note.} Preamble provenance, as required for this source: the frozen snapshot stages the chapter's own preamble file beside it (preamble.tex, 217 lines, 8116 bytes, the same shared Stacks-wide preamble the published chow and etale-cohomology cards use); the first line of the pinned chapter file is \textbackslash{}input\{preamble\}, so no preamble text lives inside the chapter itself. The card takes that file's lines 8-18 (the verbatim package and the three \textbackslash{}comment-based environments reference/slogan/history at lines 9-12, section 13.2, together with the Xy-pic setup) and its lines 154 to the runtime end of the file (the native theorem declarations and mathematical macros). Not one character was changed, no package was added, no macro was defined, and the excerpt serves only as preamble, never as a span coordinate.\par\medskip
\noindent\textit{Editorial note.} The same-chapter pointers are located at their \textbackslash{}label lines inside the pinned chapter file; the cross-chapter pointers are located in the project's own tag list staged in the snapshot (sources/stacks-chow/tags.txt), at the verbatim lines carrying their canonical tags.\par\medskip
\noindent\textit{Editorial note.} The counted proof refers to no label; the transcendental and algebraic cases are both written out inside the span, so the card carries no navigation stub and no bibliography.\par\medskip
\noindent\textit{Allowed standard prerequisites (editorial).}\begin{itemize}
\item Integrality of a multiple of a transcendental element. Used in the counted proof at lines 12160-12165. Produces an element integral over the local subring, which is what allows the extension to be kept finite.
\item Local rings inside a field and their maximal ideals. Used throughout the construction (lines 12155-12176). Ensures each enlargement is again a local subring of the same field, so that the Zorn argument applies.
\end{itemize}\medskip
\ifdefined\hypertarget\hypertarget{sourcepointer-1}{}\fi
\noindent\textit{Source reference.} Original lemma 00E0: sources/stacks-algebra/pinned-original-source.tex, source lines 3204--3204. Navigation pointer for \textbackslash{}ref\{lemma-Zariski-topology\} used inside the counted span (line 12159); the argument that consumes it is retained in the span.\par\medskip
\ifdefined\hypertarget\hypertarget{sourcepointer-2}{}\fi
\noindent\textit{Source reference.} Original lemma 00IA: sources/stacks-algebra/pinned-original-source.tex, source lines 11967--11967. Navigation pointer for \textbackslash{}ref\{lemma-dominate\} used inside the counted span (line 12162); the argument that consumes it is retained in the span.\par\medskip
\begin{lemma}
\label{lemma-find-valuation-rings}
Let $A$ be a normal domain with fraction field $K$.
\begin{enumerate}
\item For every $x \in K$, $x \not \in A$ there exists a valuation ring
$A \subset V \subset K$ with fraction field $K$ such that $x \not \in V$.
\item If $A$ is local, we can moreover choose $V$ which dominates $A$.
\end{enumerate}
In other words, $A$ is the intersection of all valuation rings in $K$
containing $A$ and if $A$ is local, then $A$ is the intersection of
all valuation rings in $K$ dominating $A$.
\end{lemma}

\begin{proof}
Suppose $x \in K$, $x \not \in A$. Consider $B = A[x^{-1}]$.
Then $x \not \in B$. Namely, if $x = a_0 + a_1x^{-1} + \ldots + a_d x^{-d}$
then $x^{d + 1} - a_0x^d - \ldots - a_d = 0$ and $x$ is integral
over $A$ in contradiction with the fact that $A$ is normal.
Thus $x^{-1}$ is not a unit in $B$. Thus $V(x^{-1}) \subset \Spec(B)$
is not empty (Lemma \ref{lemma-Zariski-topology}), and we can choose a prime
$\mathfrak p \subset B$ with $x^{-1} \in \mathfrak p$.
Choose a valuation ring $V \subset K$ dominating $B_\mathfrak p$
(Lemma \ref{lemma-dominate}).
Then $x \not \in V$ as $x^{-1} \in \mathfrak m_V$.

\medskip\noindent
If $A$ is local, then we claim that $x^{-1} B + \mathfrak m_A B \not = B$.
Namely, if $1 = (a_0 + a_1x^{-1} + \ldots + a_d x^{-d})x^{-1} +
a'_0 + \ldots + a'_d x^{-d}$ with $a_i \in A$ and $a'_i \in \mathfrak m_A$,
then we'd get
$$
(1 - a'_0) x^{d + 1} - (a_0 + a'_1) x^d - \ldots - a_d = 0
$$
Since $a'_0 \in \mathfrak m_A$ we see that $1 - a'_0$ is a unit in $A$
and we conclude that $x$ would be integral over $A$, a contradiction as
before. Then choose the prime $\mathfrak p \supset x^{-1} B + \mathfrak m_A B$
we find $V$ dominating $A$.
\end{proof}

\end{document}
